% !TeX program = pdflatex
% !TeX TXS-program:compile = txs:///pdflatex/[-shell-escape]
% !TeX TXS-program:pdflatex = pdflatex -synctex=1 -interaction=nonstopmode --shell-escape %.tex

\documentclass[../../../main_compilation.tex]{subfiles}

\begin{document}

% ================================================================================
% Encabezado / Ficha de la Máquina
% ================================================================================
\section[HTB - NOMBRE\_MAQUINA (10.10.10.X)]{NOMBRE\_MAQUINA — HackTheBox}
\label{sec:htb-nombre-maquina}

\targetSummary%
{NOMBRE\_MAQUINA}% Nombre de la máquina
{10.10.10.X}% IP
{HackTheBox}% Plataforma o imagen
{Easy}% Dificultad: Easy, Medium, Hard, Insane
{Linux}% Sistema Operativo: Linux, Windows, Android
{Web / RCE / SUID}% Categoría o Tags clave
{\today}% Fecha de resolución
{user\_flag\_hash}% User Flag
{root\_flag\_hash}% Root Flag

\vspace{0.4cm}

% ================================================================================
% 1. Resumen Ejecutivo
% ================================================================================
\subsection{Resumen Ejecutivo}
Breve descripción de la máquina, vectores de entrada y vulnerabilidades clave encontradas.

\begin{vulnbox}{Nombre de la Vulnerabilidad Principal}{High}
	\textbf{CVE:} CVE-202X-XXXX \hfill \textbf{CVSS v3.1:} 8.8 (High)\\
	\textbf{Afecta:} Nombre y versión del servicio vulnerable.\\
	\textbf{Impacto:} Resumen del impacto en el sistema.
\end{vulnbox}

% ================================================================================
% 2. Fase de Reconocimiento y Enumeración
% ================================================================================
\subsection{Fase 1: Reconocimiento y Enumeración}

\begin{terminal}[kali@pentest:\~{}]
	# 1. Escaneo rápido de puertos abiertos
	nmap -p- --open -sS --min-rate 5000 -n -Pn 10.10.10.X -oN nmap_all.txt

	# 2. Escaneo exhaustivo de versiones y scripts por defecto
	nmap -p 22,80 -sC -sV 10.10.10.X -oN nmap_services.txt
\end{terminal}

\begin{table}[h!]
	\centering
	\small
	\begin{tabularx}{\textwidth}{l l l X}
		\toprule
		\rowcolor{csecDark} \textbf{\color{white}Puerto} & \textbf{\color{white}Estado} & \textbf{\color{white}Servicio} & \textbf{\color{white}Versión / Detalles} \\
		\midrule
		\textbf{22/tcp}                                  & Abierto                      & SSH                            & OpenSSH 8.X                              \\
		\textbf{80/tcp}                                  & Abierto                      & HTTP                           & Apache 2.4.X                             \\
		\bottomrule
	\end{tabularx}
	\caption{Puertos y servicios descubiertos.}
\end{table}

\begin{notebox}[Notas de Reconocimiento]
	Anotaciones de fuzzing con gobuster/ffuf, tecnologías detectadas con Wappalyzer o subdominios descubiertos.
\end{notebox}

% ================================================================================
% 3. Fase de Explotación / Acceso Inicial
% ================================================================================
\subsection{Fase 2: Explotación (Acceso Inicial)}

\begin{codebox}[Python]{exploit.py}
	# Código / script de prueba de concepto o exploit utilizado
	import requests

	url = "http://10.10.10.X/endpoint"
	# ...
\end{codebox}

\begin{terminal}[kali@pentest:\~{}]
	nc -lvnp 4444
\end{terminal}

% ================================================================================
% 4. Escalada de Privilegios
% ================================================================================
\subsection{Fase 3: Escalada de Privilegios}

\begin{terminal}[user@target:\~{}]
	sudo -l
	# LinPEAS / WinPEAS / SUID / Tareas Cron
\end{terminal}

\begin{flagbox}{User Flag}
	user_flag_hash_aqui
\end{flagbox}

\begin{flagbox}{Root / System Flag}
	root_flag_hash_aqui
\end{flagbox}

% ================================================================================
% 5. Remediaciones
% ================================================================================
\subsection{Fase 4: Remediación y Recomendaciones}

\begin{mitigationbox}[Plan de Remediación]
	\begin{itemize}[leftmargin=*]
		\item Actualización del software a la versión corregida.
		\item Corrección de permisos y privilegios excesivos.
		\item Implementación de controles de acceso y firewall.
	\end{itemize}
\end{mitigationbox}

% ================================================================================
% 6. Lecciones y Conceptos Aprendidos
% ================================================================================
\subsection{Fase 5: Conocimientos y Técnicas Aprendidas}

\begin{learningbox}[Lecciones Asimiladas en este Reto]
	\begin{itemize}[leftmargin=*]
		\item \textbf{Vector / Técnica Aprendida:} Descripción del concepto o vector identificado.
		\item \textbf{Herramienta / Payload Clave:} Comando o payload novedoso descubierto.
		\item \textbf{Error / Obstáculo Superado:} Dificultad encontrada durante la resolución y cómo se resolvió.
	\end{itemize}
\end{learningbox}

\begin{sourcebox}[Fuentes de Consulta y Referencias del Reto]
	\begin{itemize}[leftmargin=*]
		\resourceItem{Docs}{Documentación Oficial / CVE}{Enlace o RFC correspondiente}
		\resourceItem{Libro}{Título del Libro (Autor)}{Capítulo relevante para profundizar}
		\resourceItem{Web}{Writeup / Blog de Referencia}{\url{https://...}}
	\end{itemize}
\end{sourcebox}

\end{document}
